\subsection{\problemblock{*File} Data Block}\label{sec:problem-block-file}

The \problemblock{*file} data block is used to create files containing columns of
trajectory information. A list of variable names is given in the data block, and a
column of trajectory data is written for each variable. These files are not formatted
for printing, that is, they do not contain page breaks and page headers. They are
intended for use by other applications, such as plotting programs or spreadsheets,
that require columns of data.

\problemblock{*file} data blocks can appear either within trajectory definitions or
outside of the trajectory definitions. The difference is that
\problemblock{*file} data blocks within a trajectory definition assume the
information written to the file is for a specific trajectory; whereas,
\problemblock{*file} data blocks outside of the trajectory definitions do not make
this assumption. The result is that the variables listed for
\problemblock{*file} data blocks outside of the trajectory definitions must contain
trajectory numbers enclosed in brackets.

For example,

\begin{lstlisting}[style=taosmanual]
*file compare.dat
  time[1] alt[1] alt[2] alt[3] vel[1] vel[2] vel[3]
\end{lstlisting}

writes the time from trajectory number 1, the altitudes from trajectories 1, 2, and 3,
and the velocity from trajectories 1, 2, and 3 to the file \texttt{compare.dat}.
\Cref{fig:problem-output-file-example} shows the first part of this file.

\begin{center}
  \begin{minipage}{0.96\linewidth}
    \lstinputlisting[style=taosmanual,basicstyle=\ttfamily\scriptsize]{examples/chapter04/problem-output-example.txt}
    \captionof{figure}{Problem-level trajectory output file example.}
    \label{fig:problem-output-file-example}
  \end{minipage}
\end{center}

The filename must be valid. If the file already exists, it is overwritten with no
warning and old data is lost. This is normally satisfactory because the file always
contains information from the most recent run. Any number of
\problemblock{*file} data blocks can be included in a problem; one output file is
created for each \problemblock{*file} data block.

Any standard output variable or user-defined output variable can be written to the
file. Most variables are associated with a trajectory so a trajectory number is
required, and it must be enclosed in brackets as shown above. The only exceptions
are user-defined variables defined outside of the trajectory definitions because they
are associated with the entire problem, not a single trajectory. The units and the
printout format are controlled with the \problemblock{*units/fmt} data block
(\cref{sec:block-units-fmt}).

Each output line is limited to 400 characters; this limits the number of output
variables that can be listed to about 30. Usually it is best to limit the number of
output variables to 12 or fewer because some screen editors and printers can only
handle about 130 characters per line.

Radar and relative-vehicle output variables, those starting with \texttt{rad} or
\texttt{rel}, must contain an additional subscript enclosed in square brackets. The
first subscript gives the radar station number or the trajectory number for the
relative-vehicle calculations, and the second subscript gives the vehicle trajectory
number. For example, the variable \statevar{radrng[2][1]} is used to print the radar
range of trajectory number 1 from radar station number 2, and the variable
\statevar{relvel[4][2]} is used to print the closing velocity of trajectory number 4
relative to trajectory number 2.
